\chapter{Data, Cohort Stratification \& Preprocessing}
\label{ch:data}

Three cohorts are used. DELCODE is the primary cohort on which all models are developed and
compared. ADNI and OASIS-3 were processed through an identical pipeline to serve as external
validation. This chapter documents each, the preprocessing that produced the functional
connectivity matrices, the leakage-prevention protocol, and one property of the DELCODE
follow-up schedule (Section~\ref{sec:delta-t-finding}) that turns out to constrain what can
be demonstrated on this cohort at all.

%==============================================================================
\section{The DELCODE Cohort}
\label{sec:delcode}

DELCODE (DZNE Longitudinal Cognitive Impairment and Dementia Study) is a multi-site German
observational cohort spanning the AD continuum. Diagnostic status is recorded per visit in
the \texttt{prmdiag} field, and the analysis population for this thesis is the MCI subset
with at least one usable resting-state scan.

\paragraph{Label definition.} A subject is labelled \textbf{converter} if \emph{any} visit
carries a post-conversion diagnosis, and \textbf{stable MCI} otherwise. This is a
subject-level label applied to a whole trajectory, which has two consequences that must be
handled rather than assumed away. First, the visits at which a converter is already demented
are not legitimate model inputs for a task defined as \emph{predicting} conversion; they are
excluded by the \texttt{allowed\_months} mechanism described in
Section~\ref{sec:leakage}. Second, the label is defined relative to end of follow-up, so it
is a censored observation dressed as a binary outcome --- a subject labelled stable may
simply not have converted \emph{yet}. This is the principal reason the task, as posed, is
not the clinically ideal one (Section~\ref{sec:prognoser-scope}).

\paragraph{Cohort composition.} Splits are formed once, at subject level, and reused
unchanged by every model in this thesis, so any two arms are always compared on the
same people.

\begin{table}[htbp]
  \centering
  \small
  \begin{tabular}{@{}lrrrrr@{}}
    \toprule
    \textbf{Split} & \textbf{Subjects} & \textbf{Conv.} & \textbf{Stable MCI}
      & \textbf{Mean scans} & \textbf{Max scans} \\
    \midrule
    Train      & 99  & 40 & 59 & 2.7 & 6 \\
    Validation & 34  & 14 & 20 & 2.5 & 6 \\
    Test       & 34  & 14 & 20 & 2.4 & 5 \\
    \midrule
    \textbf{Total} & \textbf{167} & \textbf{68} & \textbf{99} & 2.6 & 6 \\
    \bottomrule
  \end{tabular}
  \caption{DELCODE longitudinal MCI cohort as used throughout this thesis. Train and
    validation together form the 133-subject cross-validation pool; the 34-subject test
    split is held out and touched only for final evaluation.}
  \label{tab:delcode-cohort}
\end{table}

\risk{The held-out test set contains 34 subjects, 14 of them converters.}{%
At this size, a single subject changing side of the decision boundary moves test AUC by
roughly $0.03$--$0.04$. No difference smaller than that is measurable here, regardless of
how many decimal places a table reports, and confidence intervals on test metrics are
correspondingly wide. This is the fundamental limit on every claim in
Chapter~\ref{ch:results} and is the reason cross-validation distributions, not test point
estimates, carry the argument.}

\paragraph{The matched comparison window.} The head-to-head comparison against
Brain-TokenGT (Chapter~\ref{ch:sota-repair}) requires $2 \le T \le 3$ visits, which reduces
the cohort to $N_{\text{CV}} = 95$ (37 converters, 58 stable MCI) and $N_{\text{test}} = 25$
(11 converters, 14 stable MCI). Both models are restricted to this identical window, with
subject-identifier and fold-membership parity verified programmatically rather than inferred
from a shared random seed.

%==============================================================================
\section{The Regularity of DELCODE Follow-Up --- and Why It Matters}
\label{sec:delta-t-finding}

A time-aware longitudinal model exists to exploit irregular follow-up: to distinguish a
six-month gap from a twenty-four-month one. That premise is worth checking against the data
before claiming any benefit from it, and on DELCODE it does not hold.

\begin{table}[htbp]
  \centering
  \small
  \begin{tabular}{@{}rrr@{}}
    \toprule
    \textbf{Interval} & \textbf{Count} & \textbf{Share} \\
    \midrule
    12 months & 242 & 90.0\,\% \\
    24 months &  19 &  7.1\,\% \\
    36 months &   7 &  2.6\,\% \\
    48 months &   1 &  0.4\,\% \\
    \midrule
    \textbf{Total} & \textbf{269} & mean 13.61, SD 5.24 \\
    \bottomrule
  \end{tabular}
  \caption{Distribution of all inter-visit intervals across the 167 DELCODE subjects.
    Nine intervals in ten are identical.}
  \label{tab:delcode-intervals}
\end{table}

DELCODE is a protocol-driven study: visits are \emph{scheduled} at nominal months
$0, 12, 24, 36, 48, 60$, and subjects overwhelmingly attend on schedule. The inter-visit
interval is therefore very nearly a constant, and a model conditioned on it receives a
feature with almost no variance.

The consequence is directly measurable. Ablating the $\Delta t$ input at evaluation ---
zeroing the channel so the input shape is preserved and only the signal is removed ---
changes held-out test AUC by \textbf{exactly zero} ($0.9071$ with $\Delta t$, $0.9071$
without, identical sensitivity and specificity to four decimal places). The $\Delta t$
conditioning is architecturally present and empirically inert on this cohort. By contrast,
destroying the visit \emph{order} while preserving the $\Delta t$ distribution costs
$0.9071 \rightarrow 0.8486 \pm 0.0258$: the recurrence does use sequence order, just not
sequence timing.

\risk{The time-awareness of the model cannot be demonstrated on DELCODE, because DELCODE
has almost no time variation to be aware of.}{%
This is not a null result about the mechanism; it is a null result about the cohort. A
feature that is constant cannot contribute, and no architecture can extract information
from a variable with no variance. Any claim that $\Delta t$ conditioning helps must be
supported on a cohort where follow-up is genuinely irregular --- which is precisely what
ADNI and OASIS-3 provide (Section~\ref{sec:external-cohorts}), and which reframes external
validation from a robustness check into the only place this contribution can be tested at
all.}

\gap{Time-aware sequence models are routinely proposed and evaluated on protocol-driven
cohorts whose follow-up is regular by design.}{%
The literature on time-aware recurrence for disease progression (T-LSTM, ATTAIN, and
descendants) motivates the mechanism with clinical irregularity, but evaluation cohorts are
frequently scheduled studies in which that irregularity is largely absent. Reporting the
distribution of inter-visit intervals should be a standard prerequisite for claiming a
time-awareness benefit, and it is not.}

%==============================================================================
\section{External Validation Cohorts}
\label{sec:external-cohorts}

ADNI and OASIS-3 were carried through the same preprocessing pipeline as DELCODE so that
external validation tests the model rather than a pipeline difference.

\begin{table}[htbp]
  \centering
  \small
  \begin{tabular}{@{}lrrr@{}}
    \toprule
     & \textbf{DELCODE} & \textbf{ADNI} & \textbf{OASIS-3} \\
    \midrule
    Subjects in final manifest      & 167 & 237 & 128 \\
    Sessions in final manifest      & 434 & 567 & 234 \\
    Converter sessions              & --- & 164 & 119 \\
    Stable-MCI sessions             & --- & 403 & 106 \\
    Subjects with $\ge 2$ sessions  & 133 & 162 &  60 \\
    \quad of which converters       &  --- &  51 &  31 \\
    Visit-time encoding             & protocol month & elapsed days & elapsed days \\
    \bottomrule
  \end{tabular}
  \caption{The three cohorts after identical preprocessing and motion QC. ADNI's
    longitudinal subset is the larger of the two external cohorts and is the primary
    external validation target; OASIS-3 is retained as a second replication.}
  \label{tab:cohorts}
\end{table}

\edge{Both external cohorts encode \emph{actual elapsed days} rather than scheduled visit
labels.}{%
Where DELCODE records \texttt{\_M12\_}, ADNI and OASIS-3 record \texttt{ses-d0381}. Their
follow-up is genuinely irregular, which makes them the only cohorts in this project on
which the $\Delta t$ mechanism of Section~\ref{sec:delta-t-finding} can be shown to
contribute anything. Combined with a five-fold larger longitudinal converter count than
DELCODE offers, this is the single highest-value asset the project holds.}

\risk{Naively reusing the DELCODE visit parser on these cohorts fails silently.}{%
A month parser written against \texttt{\_M\textit{N}\_} filename tokens returns
\texttt{None} for every ADNI and OASIS-3 session; the loader then skips those visits, each
subject lands with zero visits, and the cohort disappears without raising. The result is a
clean, well-formatted, entirely empty evaluation. The cohort-aware parsing described in
Section~\ref{sec:visit-semantics} exists specifically to make this failure loud instead of
silent.}

%==============================================================================
\section{fMRI Preprocessing \& Confound Management}
\label{sec:preprocessing}

All three cohorts pass through the same sequence: BIDS organisation, fMRIPrep, denoising,
motion-QC gating, atlas extraction, and connectivity estimation.

\begin{enumerate}
  \item \textbf{BIDS conversion.} DICOM archives (ADNI) and raw NIfTI pools (OASIS-3) are
    converted to BIDS. Multi-run acquisitions are concatenated --- 6 series merged in ADNI,
    and 587 raw runs consolidated into 284 single-session 4-D volumes in OASIS-3, where
    short scout runs and repeated rest runs are common.
  \item \textbf{fMRIPrep.} Standard motion correction, susceptibility distortion
    correction, coregistration to the subject's T1w anatomical, and normalisation. A T1w
    scan is required, which is itself an attrition gate.
  \item \textbf{Denoising.} ICA-AROMA plus two physiological regressors and one global
    signal regressor (\texttt{ICAAROMA2Phys1GS}), followed by bandpass filtering.
  \item \textbf{Motion QC gate.} Sessions with mean framewise displacement above
    $0.5$\,mm are excluded. This is applied identically across cohorts.
  \item \textbf{Atlas extraction and connectivity.} Schaefer-200 parcel time series are
    extracted, and pairwise Pearson correlation followed by the Fisher $z$-transform yields
    the $200 \times 200$ matrix per session.
  \item \textbf{Age-effect regression.} Linear-model regression of age effects, so that
    normal ageing is not learned as disease.
\end{enumerate}

\paragraph{Attrition is fully accounted.} Of the 414 subjects completing fMRIPrep and
denoising across the two external cohorts, 396 reach the final flat product. Every one of
the 18 missing subjects has a documented, per-subject cause: 16 (12 OASIS-3, 4 ADNI) had
\emph{all} sessions fail the motion gate, and 2 OASIS-3 subjects
(\texttt{sub-OAS30797}, \texttt{sub-OAS31416}) lack an fMRIPrep confounds file, so QC could
not be evaluated at all. At the session level the motion gate removes 40 of 714 ADNI
sessions (5.6\%) and 36 of 270 OASIS-3 sessions (13.3\%); excluded scans average
$>0.67$\,mm mean FD with over half of their frames exceeding $0.5$\,mm displacement.

\edge{Counting artifacts of the filesystem is not the same as counting data --- a lesson
learned twice.}{%
Two separate pipeline-monitoring bugs in this project inflated apparent completion by
counting \emph{directory existence} rather than expected file content. In the more severe
instance, 46 OASIS-3 subjects had an fMRIPrep output directory and \emph{zero} postprocessed
volumes, so a progress monitor reported 90\% settled while 36\% of what it counted was
empty. All 46 proved recoverable, and their recovery moved the OASIS-3 longitudinal cohort
from 40 to 60 subjects (19 to 31 converters). Every count reported in this thesis is
therefore defined as \emph{files matching the expected content pattern}, never as a
directory count.}

%==============================================================================
\section{Cross-Cohort Visit Semantics}
\label{sec:visit-semantics}

Section~\ref{sec:gap-cohorts} identified the conflation of three distinct quantities as a
silent failure mode. They are kept separate throughout, with one function per concept
dispatching on cohort:

\begin{description}[leftmargin=3.2cm,style=nextline]
  \item[\texttt{visit\_index}] Zero-based rank of the session within the subject's
    trajectory. Used for ordering and window selection. Always available.
  \item[\texttt{protocol\_month}] The nominal scheduled-visit label. Used \emph{only} by the
    \texttt{allowed\_months} leakage filter. Returns \texttt{None} when a session cannot be
    confidently mapped onto a scheduled visit --- for example an ADNI unscheduled visit code
    --- and is never guessed.
  \item[\texttt{delta\_t\_months}] Cumulative elapsed months since baseline, as a float.
    This is the model input. Strictly increasing per subject by construction, so a single
    inter-visit-difference code path serves all three cohorts.
\end{description}

For DELCODE, \texttt{delta\_t\_months} equals the nominal protocol month exactly, so the
cohort-aware refactor is required to reproduce existing DELCODE results byte-for-byte. That
equality is used as an acceptance gate rather than as an expectation.

%==============================================================================
\section{Methodological Rigour \& Data Leakage Prevention}
\label{sec:leakage}

Section~\ref{sec:gap-floors} noted that connectome studies are systematically inflated by
leakage, with scan-level cross-validation alone documented to add 30--55 accuracy points.
Four mechanisms guard against it here.

\paragraph{Patient-level splitting.} Splits are formed over \emph{subjects}, never over
scans. A subject's every visit lies wholly within one split. Cross-validation uses stratified
grouped folds with the subject identifier as the group, so the same guarantee holds
fold-to-fold.

\paragraph{An isolated test manifest.} The 34-subject test split is written once and read
only at final evaluation. It participates in no model selection, no threshold selection, and
no early stopping.

\paragraph{Post-conversion visit exclusion.} For converter subjects, visits already
relabelled as dementia are removed from the \texttt{allowed\_months} list. Without this, the
classifier would be shown scans of established dementia and asked to predict conversion ---
a leak of the label into the input.

\paragraph{Thresholds derived from validation only.} Decision thresholds come from
out-of-fold validation predictions, never from test metrics. This is enforced in code rather
than by convention: the evaluation entry point raises rather than defaulting when no explicit
threshold is supplied (Section~\ref{sec:thresholds}).

\paragraph{An executable audit.} Split hygiene is not asserted in prose. A
\texttt{run\_full\_audit} call at the head of every training notebook hard-fails if any
subject identifier appears in more than one split, and a dedicated experiment
(\texttt{sanity-split-hygiene}) exists solely to run that check. Appendix~B reproduces the
resulting logs as evidence rather than summarising them.

\edge{Leakage prevention here is executable and fails loudly, rather than being described.}{%
The distinction matters because leakage is not usually introduced deliberately; it is
introduced by a refactor that nobody re-audited. A check that runs on every training job and
aborts the run is the only form of this guarantee that survives six months of development.}
